Para cerrar, nos quedamos con tres lecciones clave de toda la exposición.
Primero, ningún protocolo es mejor en abstracto: TCP prioriza entrega confiable, UDP evita el handshake y las retransmisiones, y WebSockets resuelve la comunicación bidireccional continua.
Segundo, toda decisión de transporte implica un compromiso: al quitar controles de red reducimos sobrecarga, pero la aplicación debe tolerar o gestionar posibles pérdidas.
Y tercero, en proyectos reales estos protocolos no compiten, se combinan según lo que necesita cada módulo.
En definitiva, elegir el protocolo correcto desde el diseño es lo que permite que un sistema distribuido funcione bien en producción.